\chapter{Keamanan Aplikasi Web}
\label{ch:keamanan}

\section*{Tujuan Pembelajaran}
\begin{itemize}
  \item Menyusun model ancaman satu aplikasi web, yaitu menyebut batas
        kepercayaannya, data yang bernilai di dalamnya, dan jalan masuk yang
        tersedia bagi penyerang.
  \item Memasang pertahanan bawaan pada sisi server, yaitu header keamanan,
        atribut cookie, daftar origin, \textit{query} berparameter, dan
        \textit{escaping} keluaran, lalu membuktikan terpasangnya dari
        \textit{response} yang dikirim aplikasi.
  \item Mengukur keamanan sebagai angka, yaitu jumlah syarat yang terpenuhi
        dan jumlah kerentanan dependensi, lalu menegakkannya sebagai
        \textit{gate} yang memblokir \textit{merge}.
\end{itemize}

\section{Pendahuluan}

Lima bab sebelumnya mengejar cepat dan benar. Bab ini menambahkan satu
pertanyaan lagi: apa yang terjadi bila ada yang memakai aplikasi ini dengan
cara yang tidak diniatkan pembuatnya. Pertanyaan itu tidak dapat dijawab
dengan mencoba-coba, karena jumlah cara memakai sebuah aplikasi jauh lebih
banyak daripada jumlah cara memakainya dengan benar.

Jawabannya disusun terbalik. Alih-alih menebak satu per satu serangan yang
mungkin, pertahanannya dipasang sebagai bawaan, lalu keberadaannya diperiksa
ulang pada tiap perubahan. Cara ini memindahkan keamanan dari ingatan orang
ke dalam file yang dapat dibaca mesin, persis seperti \textit{budget} performa
pada Bab 5 dan batas mutu pada Bab 6.

Seluruh pemeriksaan pada bab ini bersifat mengamati dan dijalankan terhadap
aplikasi sendiri di komputer sendiri. Referensi utamanya adalah daftar risiko
Open Worldwide Application Security Project (OWASP) Top 10
\cite{owasp-top10} untuk memilih apa yang dijaga, Content
Security Policy Level 3 \cite{w3c-csp3} untuk header, RFC 6265
\cite{rfc6265} untuk cookie, dan Fetch Standard \cite{whatwg-fetch} untuk
aturan lintas origin.

\section{Threat Modelling}

\textit{Threat modelling} adalah menuliskan lebih dahulu apa yang dijaga, dari
siapa, dan lewat jalan mana \cite{owasp-threat-modeling}. Hasilnya bukan
daftar serangan, melainkan daftar batas: di mana data berpindah tangan, dan
siapa yang dipercaya di tiap sisinya.

Analoginya seperti menata ruang arsip fakultas. Yang ditulis lebih dahulu
bukan daftar cara orang menyelinap, melainkan daftar pintu, siapa yang
memegang kuncinya, dan lemari mana yang berisi berkas bernilai. Sesudah
daftar itu ada, pengamanannya tinggal dipasang di tiap pintu.

Gambar \ref{fig:batas-percaya} memperlihatkan batas kepercayaan aplikasi bab
ini.

\begin{figure}[!htb]
  \centering
  \begin{tikzpicture}[
    font=\scriptsize,
    kotak/.style={draw, rounded corners, minimum height=0.9cm,
                  minimum width=2.3cm, align=center, font=\scriptsize\bfseries},
    panah/.style={-{Stealth[length=2mm]}, thick}
  ]
    \node[kotak, fill=red!22] (browser) at (0,0) {Browser\\{\tiny\mdseries tidak dipercaya}};
    \node[kotak, fill=orange!12] (server) at (4.2,0) {Aplikasi\\{\tiny\mdseries memeriksa semua}};
    \node[kotak, fill=blue!8] (basis) at (8.4,0) {Basis data\\{\tiny\mdseries di jaringan sendiri}};
    \node[kotak, fill=gray!35] (paket) at (4.2,-1.8) {Paket pihak lain\\{\tiny\mdseries ikut berjalan}};

    \draw[panah] (browser) -- node[font=\tiny, above] {teks pengguna} (server);
    \draw[panah] (server) -- node[font=\tiny, above] {query berparameter} (basis);
    \draw[panah] (paket) -- (server);
    \draw[dashed, thick, gray] (2.1,1.1) -- (2.1,-2.6);
    \node[font=\tiny, gray, anchor=south] at (2.1,1.15) {batas kepercayaan};
  \end{tikzpicture}
  \caption{Semua yang datang dari kiri garis putus-putus diperlakukan sebagai
           teks yang belum dipercaya, termasuk paket pihak lain yang ikut
           berjalan di dalam aplikasi}
  \label{fig:batas-percaya}
\end{figure}

Dari gambar tersebut keluar tiga kewajiban yang dipakai sepanjang bab ini.
Pertama, teks dari browser tidak pernah disambung langsung ke dalam perintah
apa pun. Kedua, \textit{response} membawa aturan main untuk browser, yaitu
header keamanan dan atribut cookie. Ketiga, paket pihak lain ikut dihitung
sebagai jalan masuk, sehingga daftarnya diperiksa berkala.

\section{Header Keamanan}

Header keamanan adalah aturan yang dikirim server untuk dipatuhi browser.
Header itu tidak mengubah apa pun di sisi server, dan justru itulah
kekuatannya: satu baris pada \textit{response} mempersempit apa yang boleh
dilakukan halaman di komputer pengunjung.

Tabel \ref{tab:header} membandingkan \textit{response} aplikasi bab ini
dengan \textit{response} server statis bawaan Python yang tidak disetel sama
sekali. Keduanya diukur dengan \texttt{curl} pada komputer yang sama.

\begin{table}[!htb]
  \centering
  \caption{Header keamanan pada server bawaan lawan aplikasi yang sudah
           dikeraskan, hasil pengukuran nyata}
  \label{tab:header}
  \begin{tabularx}{\textwidth}{Xll}
    \hline
    \textbf{Header} & \textbf{Server bawaan} & \textbf{Aplikasi Bab 7} \\
    \hline
    \texttt{Content-Security-Policy} & tidak ada & \texttt{default-src 'self'} \\
    \hline
    \texttt{X-Content-Type-Options} & tidak ada & \texttt{nosniff} \\
    \hline
    \texttt{Referrer-Policy} & tidak ada & \texttt{no-referrer} \\
    \hline
    \texttt{Permissions-Policy} & tidak ada & \texttt{geolocation=(), camera=(), microphone=()} \\
    \hline
  \end{tabularx}
\end{table}

Empat baris itu menjaga empat hal yang berbeda. \texttt{Content-Security-Policy}
membatasi dari mana skrip dan gambar boleh dimuat, sehingga kode yang
berhasil diselipkan ke halaman tetap tidak punya tempat mengambil
perintahnya \cite{w3c-csp3}. \texttt{X-Content-Type-Options} menutup
kebiasaan browser menebak jenis file dari isinya. \texttt{Referrer-Policy}
menahan alamat halaman asal agar tidak ikut bocor ke situs lain.
\texttt{Permissions-Policy} mematikan kemampuan perangkat yang memang tidak
dipakai, yaitu lokasi, kamera, dan mikrofon.

Listing \ref{lst:header} memperlihatkan tempat keempatnya dipasang. Letaknya
di satu fungsi yang dilewati seluruh \textit{response}, sehingga alamat baru
tidak dapat lupa memasangnya.

\begin{lstlisting}[style=PythonStyle, caption={Header keamanan dipasang di satu tempat, tersedia di \berkas{sources/chapter07/aplikasi.py}}, label={lst:header}]
  def kirim_response(self, status, isi_response, jenis_isi, header_tambahan):
    """Mengirim satu response lengkap beserta seluruh header keamanannya."""
    self.send_response(status)
    self.send_header("Content-Type", jenis_isi)
    self.send_header("Content-Length", str(len(isi_response)))
    for nama_header, nilai in self.server.kebijakan["header_wajib"].items():
      self.send_header(nama_header, nilai)
\end{lstlisting}

\section{Cookie yang Aman}

Cookie sesi adalah penanda yang membuat server mengenali pengunjung yang sama
pada permintaan berikutnya. Karena isinya cukup untuk mewakili seseorang,
cookie diberi atribut pembatas, bukan dikirim polos \cite{rfc6265}.

Tabel \ref{tab:cookie} memuat atribut yang dipakai aplikasi bab ini beserta
yang dijaga masing-masing. Nilainya diambil apa adanya dari header
\texttt{Set-Cookie} yang terbaca \texttt{curl}.

\begin{table}[!htb]
  \centering
  \caption{Atribut cookie sesi aplikasi Bab 7, hasil pembacaan header
           \texttt{Set-Cookie}}
  \label{tab:cookie}
  \begin{tabularx}{\textwidth}{lX}
    \hline
    \textbf{Atribut} & \textbf{Yang dijaga} \\
    \hline
    \texttt{HttpOnly} & Cookie tidak dapat dibaca JavaScript mana pun yang berjalan di halaman \\
    \hline
    \texttt{SameSite=Lax} & Cookie tidak ikut terkirim pada permintaan yang dipicu situs lain \\
    \hline
    \texttt{Path=/} & Jangkauan cookie ditulis eksplisit, bukan ditebak browser \\
    \hline
    \texttt{Max-Age=3600} & Sesi berakhir sendiri sesudah satu jam \\
    \hline
  \end{tabularx}
\end{table}

Satu atribut sengaja tidak dipasang, yaitu \texttt{Secure}. Atribut itu
membuat cookie hanya dikirim lewat HTTPS, sedangkan latihan bab ini berjalan
lewat HTTP di \textit{localhost}. Pada aplikasi yang sungguhan dilayani lewat
HTTPS, atribut itu wajib ada, dan alasannya ditulis di komentar kode agar
tidak hilang ketika aplikasinya dipindahkan ke server.

\texttt{SameSite=Lax} juga menjawab satu risiko yang sering disebut
tersendiri, yaitu permintaan yang dipicu dari situs lain memakai sesi
pengunjung yang sedang masuk. Dengan atribut itu, cookie tidak ikut terkirim
pada permintaan semacam itu, sehingga server memperlakukannya sebagai
pengunjung yang belum dikenal.

\section{Aturan Lintas Origin}

\textit{Cross-Origin Resource Sharing} (CORS) sering disalahpahami sebagai
pengaman server, padahal yang dilindunginya adalah pembaca
\cite{whatwg-fetch}. Browser menahan halaman dari satu situs untuk membaca
\textit{response} milik situs lain, kecuali situs pemilik \textit{response}
itu menyatakan izin lewat header.

Dua hal pada CORS sering tertukar. Pembagiannya begini. Izin CORS menentukan
siapa yang boleh membaca \textit{response}, sedangkan pemeriksaan hak akses
menentukan siapa yang boleh meminta. Server tetap wajib memeriksa hak akses
sendiri, karena permintaan dari luar browser tidak terikat aturan CORS sama
sekali.

Tabel \ref{tab:cors} memuat perilaku aplikasi bab ini untuk origin yang
terdaftar dan yang tidak.

\begin{table}[!htb]
  \centering
  \caption{Izin lintas origin menurut daftar di
           \berkas{sources/chapter07/kebijakan.json}, hasil pengukuran nyata}
  \label{tab:cors}
  \begin{tabularx}{\textwidth}{lX}
    \hline
    \textbf{Origin yang meminta} & \textbf{Header izin yang dikirim} \\
    \hline
    \texttt{https://toko.pradita.ac.id} & \texttt{Access-Control-Allow-Origin} berisi origin itu sendiri \\
    \hline
    \texttt{http://127.0.0.1:8030} & \texttt{Access-Control-Allow-Origin} berisi origin itu sendiri \\
    \hline
    \texttt{https://situs-lain.example} & tidak ada, hanya \texttt{Vary: Origin} \\
    \hline
  \end{tabularx}
\end{table}

Nilai bintang sengaja tidak dipakai. Bintang membuka \textit{response} bagi
situs mana pun, dan sekali dipasang, menariknya kembali berarti merusak
halaman yang sudah terlanjur bergantung padanya. Header \texttt{Vary: Origin}
ikut dikirim agar \textit{cache} perantara tidak menyimpan satu izin lalu
memakainya untuk origin yang lain.

\section{Teks dari Pengguna}

Dua risiko terbesar pada daftar OWASP lahir dari kebiasaan yang sama, yaitu
menyambung teks pengguna ke dalam sesuatu yang kemudian dijalankan
\cite{owasp-top10}. Jawabannya pun sama bentuknya: teks itu dikirim sebagai
data, bukan sebagai bagian perintah.

Pada basis data, caranya memakai parameter. Listing \ref{lst:parameter}
memperlihatkan bentuknya. Teks apa pun yang dikirim pengunjung berakhir
sebagai satu nilai pencarian, tidak peduli tanda baca apa yang ada di
dalamnya.

\begin{lstlisting}[style=PythonStyle, caption={Teks pengguna dikirim sebagai parameter, bukan disambung ke SQL, tersedia di \berkas{sources/chapter07/aplikasi.py}}, label={lst:parameter}]
def cari_buku(judul_dicari):
  """Mencari buku menurut judulnya, dengan judul dikirim sebagai parameter.

  Teks dari pengguna tidak pernah disambung ke dalam SQL. Tanda kutip, tanda
  titik koma, dan kata kunci SQL di dalamnya diperlakukan sebagai data biasa,
  sehingga hasilnya hanya kosong, bukan perintah tambahan.
  """
  pola = f"%{judul_dicari}%"
  with psycopg.connect(DSN) as koneksi:
    baris = koneksi.execute(SQL_CARI_BUKU, (pola, BATAS_HASIL_CARI)).fetchall()
\end{lstlisting}

Tabel \ref{tab:parameter} memuat dua pencarian yang dijalankan terhadap
aplikasi bab ini. Baris kedua memakai kata kunci yang di dalamnya ada tanda
kutip dan tanda titik koma.

\begin{table}[!htb]
  \centering
  \caption{Pencarian dengan teks biasa dan dengan teks bertanda baca, hasil
           pengukuran nyata}
  \label{tab:parameter}
  \begin{tabularx}{\textwidth}{Xrl}
    \hline
    \textbf{Kata kunci yang dikirim} & \textbf{Baris} & \textbf{Status} \\
    \hline
    \texttt{Pengantar Jilid 7} & 20 & 200, tanpa \textit{error} \\
    \hline
    \texttt{Pengantar'; SELECT 1 -{}-} & 0 & 200, tanpa \textit{error} \\
    \hline
  \end{tabularx}
\end{table}

Baris kedua itu yang perlu dibaca pelan. Hasilnya nol bukan karena teksnya
ditolak, melainkan karena tidak ada judul buku yang berisi rangkaian huruf
tersebut. Bagi basis data, isinya tidak berbeda dengan judul biasa, dan di
situlah letak pengamanannya.

Pada HTML, bentuk jawabannya adalah \textit{escaping} keluaran. Teks pengguna
diubah lebih dahulu sehingga tanda kurung siku dan tanda kutip terbaca
sebagai huruf, bukan sebagai markup. Pengujian pada aplikasi ini mengirim
\texttt{<b>Mahasiswa</b>} sebagai nama, dan halaman yang keluar berisi
\texttt{\&lt;b\&gt;Mahasiswa\&lt;/b\&gt;}, yaitu teks yang tampil apa adanya
di layar.

\section{Manajemen Secret}

Secret adalah nilai yang kekuatannya habis begitu diketahui orang lain,
misalnya kunci penanda sesi, kata sandi basis data, dan token layanan luar.
Aturannya satu kalimat: secret tidak pernah ditulis di dalam kode, karena
kode masuk repositori dan dibaca siapa pun yang punya aksesnya.

Listing \ref{lst:secret} memperlihatkan cara aplikasi bab ini membacanya.

\begin{lstlisting}[style=PythonStyle, caption={Secret dibaca dari environment, bukan dari kode, tersedia di \berkas{sources/chapter07/aplikasi.py}}, label={lst:secret}]
def ambil_secret_sesi():
  """Mengambil secret penanda sesi dari environment.

  Secret tidak pernah ditulis di kode, karena kode masuk repositori dan
  dibaca siapa saja yang punya aksesnya. Bila environment-nya kosong,
  aplikasi membuat nilai acak sekali jalan, yang hilang begitu aplikasi
  dimatikan.
  """
  return os.environ.get("SECRET_SESI") or secrets.token_hex(PANJANG_SESI_BYTE)
\end{lstlisting}

Tabel \ref{tab:secret} merangkum tempat tiap jenis nilai disimpan. Yang
membedakan bukan tingkat kerahasiaannya saja, melainkan juga apa yang harus
dilakukan ketika nilainya bocor.

\begin{table}[!htb]
  \centering
  \caption{Tempat menyimpan nilai menurut jenisnya}
  \label{tab:secret}
  \begin{tabularx}{\textwidth}{lXl}
    \hline
    \textbf{Jenis nilai} & \textbf{Contoh} & \textbf{Tempatnya} \\
    \hline
    Setelan biasa & Port, nama basis data, batas halaman & File kode \\
    \hline
    Setelan per lingkungan & Alamat basis data, daftar origin & File kebijakan \\
    \hline
    Secret & Kunci sesi, kata sandi, token & Environment atau \textit{secret manager} \\
    \hline
  \end{tabularx}
\end{table}

\section{Keamanan Rantai Pasok}

Paket pihak lain ikut berjalan di dalam aplikasi dengan hak yang sama.
Karena itu daftar dependensi diperlakukan sebagai bagian dari permukaan
serangan, dan diperiksa berkala terhadap pengumuman kerentanan
\cite{pypa-advisory}.

Tabel \ref{tab:dependensi} memuat hasil pemeriksaan pada repositori modul ini.

\begin{table}[!htb]
  \centering
  \caption{Hasil pemeriksaan dependensi, hasil pengukuran nyata}
  \label{tab:dependensi}
  \begin{tabularx}{\textwidth}{lrX}
    \hline
    \textbf{Ekosistem} & \textbf{Paket} & \textbf{Temuan} \\
    \hline
    Python, lewat \texttt{pip-audit} & 48 & Tidak ada kerentanan yang diumumkan \\
    \hline
    Node, lewat \texttt{npm audit} & paket Bab 6 & Tidak ada kerentanan yang diumumkan \\
    \hline
  \end{tabularx}
\end{table}

Angka nol pada tabel tersebut punya umur simpan. Pemeriksaan ini mencocokkan
versi yang terpasang dengan pengumuman yang sudah ada hari itu, sehingga
paket yang hari ini bersih dapat berubah status besok tanpa satu baris kode
pun berubah. Karena itu pemeriksaannya dipasang berulang di
\textit{Continuous Integration}, bukan dijalankan sekali saat memasang.

\section{Gate Keamanan di CI}

Seluruh pertahanan di atas mudah terpasang dan sama mudahnya hilang, misalnya
ketika sebuah alamat baru ditulis tanpa melewati fungsi pengirim
\textit{response}. Karena itu syaratnya ditulis sebagai file, lalu dinilai
ulang pada tiap perubahan. Gambar \ref{fig:gate-keamanan} memperlihatkan
tempatnya bekerja.

\begin{figure}[!htb]
  \centering
  \begin{tikzpicture}[
    font=\scriptsize,
    kotak/.style={draw, rounded corners, minimum height=0.9cm, minimum width=2.1cm,
                  align=center, font=\scriptsize\bfseries},
    panah/.style={-{Stealth[length=2mm]}, thick}
  ]
    \node[kotak, fill=blue!8] (ubah) at (0,0) {Perubahan\\kode};
    \node[kotak, fill=blue!8] (baca) at (3.0,0) {Baca response\\{\tiny\mdseries header, cookie, CORS}};
    \node[kotak, fill=orange!12] (nilai) at (6.2,0) {Nilai syarat\\{\tiny\mdseries kebijakan.json}};
    \node[kotak, fill=green!10] (gabung) at (9.6,0.9) {Status 0\\{\tiny\mdseries boleh digabung}};
    \node[kotak, fill=red!22] (tolak) at (9.6,-0.9) {Status 1\\{\tiny\mdseries merge diblokir}};

    \draw[panah] (ubah) -- (baca);
    \draw[panah] (baca) -- (nilai);
    \draw[panah] (nilai.east) -- ++(0.5,0) |- (gabung.west);
    \draw[panah] (nilai.east) -- ++(0.5,0) |- (tolak.west);
    \node[font=\tiny, gray, anchor=east] at (7.7,0.72) {seluruh syarat terpenuhi};
    \node[font=\tiny, gray, anchor=east] at (7.7,-0.72) {ada syarat yang hilang};
  \end{tikzpicture}
  \caption{Syarat keamanan dinilai dari \textit{response} yang benar-benar
           dikirim aplikasi, bukan dari pembacaan kode}
  \label{fig:gate-keamanan}
\end{figure}

Listing \ref{lst:gate-keamanan} memuat keluaran gerbangnya pada aplikasi bab
ini, dipersingkat pada bagian yang berulang.

\begin{lstlisting}[language=bash, caption={Menjalankan gate keamanan, tersedia di \berkas{sources/chapter07/periksa-keamanan.py}}, label={lst:gate-keamanan}]
source ../.venv/bin/activate
python periksa-keamanan.py
echo $?

# Keluarannya:
#   Pemeriksaan                                  Keadaan      Verdict
#   / Content-Security-Policy                    sesuai       lolos
#   / X-Content-Type-Options                     sesuai       lolos
#   / Referrer-Policy                            sesuai       lolos
#   / Permissions-Policy                         sesuai       lolos
#   cookie sesi HttpOnly                         sesuai       lolos
#   cookie sesi SameSite=Lax                     sesuai       lolos
#   CORS izin untuk https://toko.pradita.ac.id   sesuai       lolos
#   CORS tolak https://situs-lain.example        sesuai       lolos
#
# Lolos: seluruh syarat keamanan terpenuhi.
# 0
\end{lstlisting}

Gerbang yang selalu lolos tidak membuktikan apa pun, sehingga pemeriksaan
yang sama dijalankan terhadap pembanding, yaitu server statis bawaan Python
tanpa penyetelan. Tabel \ref{tab:gate-banding} memuat kedua hasilnya.

\begin{table}[!htb]
  \centering
  \caption{Verdict gerbang pada dua sasaran, hasil pengukuran nyata}
  \label{tab:gate-banding}
  \begin{tabularx}{\textwidth}{Xrrr}
    \hline
    \textbf{Sasaran} & \textbf{Syarat} & \textbf{Lolos} & \textbf{Status keluar} \\
    \hline
    Aplikasi Bab 7 & 18 & 18 & 0 \\
    \hline
    Server statis bawaan Python & 9 & 1 & 1 \\
    \hline
  \end{tabularx}
\end{table}

Satu baris pada pembanding justru lolos, dan sebabnya layak diperhatikan.
Syarat itu berbunyi bahwa origin yang tidak terdaftar tidak boleh menerima
izin lintas origin, sedangkan server bawaan memang tidak pernah mengirim
header CORS sama sekali. Syarat yang ditulis sebagai larangan dapat terpenuhi
karena fiturnya tidak ada, bukan karena pengamanannya bekerja, dan syarat
semacam itu perlu dipasangkan dengan syarat positifnya.

Analisis statis melengkapi gerbang tersebut dari sisi kode. Pengaturan
\texttt{ruff} pada bab ini menyalakan aturan keamanan dari bandit, misalnya
peringatan untuk SQL yang disambung dari potongan teks dan untuk nilai yang
tampak seperti kata sandi tertulis di kode. Pada seluruh file sumber Bab 7,
pemeriksaan itu melaporkan nol temuan.

\section{Ringkasan}

Keamanan tidak dikejar dengan menebak satu per satu cara orang menyalahgunakan
aplikasi, melainkan dengan menuliskan apa yang dijaga dan dari mana jalan
masuknya. Model ancaman bab ini berisi satu batas kepercayaan, dan dari batas
itu keluar tiga kewajiban: teks pengguna tidak pernah disambung ke perintah,
\textit{response} membawa aturan main untuk browser, dan paket pihak lain
dihitung sebagai jalan masuk.

Pertahanan sisi server dipasang sebagai bawaan, bukan pilihan. Empat header
keamanan dikirim pada tiap \textit{response} lewat satu fungsi, cookie sesi
diberi \texttt{HttpOnly} dan \texttt{SameSite=Lax}, dan izin lintas origin
hanya diberikan kepada daftar yang ditulis di file kebijakan. Server bawaan
Python yang dipakai sebagai pembanding tidak mengirim satu pun dari keempat
header itu.

Teks dari pengguna diperlakukan sebagai data di mana pun ia berakhir.
Pencarian dengan kata kunci bertanda kutip dan titik koma dijawab 200 dengan
nol baris, tanpa \textit{error}, karena isinya hanya dibandingkan sebagai
teks. Nama yang berisi markup tampil di layar sebagai huruf, karena
keluarannya di-\textit{escape} lebih dahulu.

Nilai yang kekuatannya habis begitu diketahui orang lain disimpan di luar
kode, dan daftar dependensi diperiksa terhadap pengumuman kerentanan. Empat
puluh delapan paket Python dan paket Node Bab 6 sama-sama bersih pada
pemeriksaan hari itu, dan justru karena jawabannya bergantung pada hari,
pemeriksaannya diulang di \textit{Continuous Integration}.

Seluruhnya berakhir pada satu angka yang dapat dibaca mesin. Gerbang keamanan
membaca \textit{response} yang benar-benar dikirim, menilainya terhadap 18
syarat, lalu keluar dengan status 0. Pemeriksaan yang sama terhadap server
tanpa penyetelan meloloskan satu syarat dari sembilan dan keluar dengan
status 1. Selisih itulah bukti bahwa yang dinilai memang pertahanannya, bukan
niat baik penulisnya.

\section{Latihan}

Seluruh file yang dipakai pada latihan ini tersedia di
\berkas{sources/chapter07/}, dihitung relatif terhadap akar repositori.
Latihan membutuhkan Docker dengan Compose versi 2 dan Python 3. Listing
penyiapan masuk ke direktori itu lebih dahulu, dan seluruh perintah sesudahnya
dijalankan dari sana.

Seluruh pemeriksaan pada latihan ini bersifat mengamati dan hanya diarahkan
ke aplikasi yang berjalan di komputer sendiri. Tidak ada satu pun perintah di
bab ini yang boleh diarahkan ke situs kampus maupun sistem milik orang lain.
Hasil tiap mahasiswa akan berbeda pada bagian waktu dan daftar paket, dan
yang dilaporkan adalah pola perbandingannya.

Penyiapannya dilakukan sekali, seperti pada Listing \ref{lst:c7-siapkan}.

\begin{lstlisting}[language=bash, caption={Menyiapkan basis data dan venv, memakai \berkas{sources/chapter07/siapkan-data.sh}}, label={lst:c7-siapkan}]
cd sources/chapter07
docker compose up -d
./siapkan-data.sh
python3 -m venv ../.venv
../.venv/bin/pip install "psycopg[binary]" ruff pip-audit
\end{lstlisting}

\subsection*{Latihan 1: Membaca Header Keamanan}

Latihan ini membuktikan bahwa header keamanan adalah pilihan yang harus
dibuat, bukan sesuatu yang datang sendiri.

\begin{enumerate}
  \item Nyalakan aplikasinya, lalu nyalakan juga server statis bawaan Python
        sebagai pembanding, seperti pada Listing \ref{lst:c7-lat1-jalan}.
  \item Baca header kedua server dengan \texttt{curl}, lalu isi Tabel
        \ref{tab:c7-lat1-hasil}. Tabel \ref{tab:c7-lat1-contoh}
        memperlihatkan contoh yang sudah terisi.
  \item Tambahkan satu header baru ke \berkas{sources/chapter07/kebijakan.json},
        misalnya \texttt{X-Frame-Options} bernilai \texttt{DENY}, lalu
        nyalakan ulang aplikasinya dan periksa kembali.
  \item Jawab tiga pertanyaan berikut. Pertama, header mana yang membatasi
        asal skrip yang boleh dimuat halaman? Kedua, mengapa header keamanan
        dipasang di satu fungsi dan bukan di tiap alamat? Ketiga, apa yang
        terjadi pada header baru tadi tanpa aplikasinya dinyalakan ulang?
\end{enumerate}

\begin{lstlisting}[language=bash, caption={Membaca header dua server, tersedia di \berkas{sources/chapter07/aplikasi.py}}, label={lst:c7-lat1-jalan}]
source ../.venv/bin/activate
SECRET_SESI=rahasia-latihan python aplikasi.py &
(cd statis && python -m http.server 8031 &)
curl -sI http://127.0.0.1:8030/ | grep -iE "^(content-security|x-content|referrer|permissions)"
curl -sI http://127.0.0.1:8031/index.html | grep -iE "^(content-security|x-content|referrer|permissions)"

# Keluarannya:
# Content-Security-Policy: default-src 'self'
# X-Content-Type-Options: nosniff
# Referrer-Policy: no-referrer
# Permissions-Policy: geolocation=(), camera=(), microphone=()
\end{lstlisting}

\begin{table}[!htb]
  \centering
  \caption{Contoh hasil Latihan 1 yang sudah terisi}
  \label{tab:c7-lat1-contoh}
  \begin{tabularx}{\textwidth}{Xll}
    \hline
    \textbf{Header} & \textbf{Server bawaan} & \textbf{Aplikasi Bab 7} \\
    \hline
    \texttt{Content-Security-Policy} & tidak ada & \texttt{default-src 'self'} \\
    \hline
    \texttt{X-Content-Type-Options} & tidak ada & \texttt{nosniff} \\
    \hline
  \end{tabularx}
\end{table}

\begin{table}[!htb]
  \centering
  \caption{Lembar isian hasil pengamatan Latihan 1}
  \label{tab:c7-lat1-hasil}
  \begin{tabularx}{\textwidth}{Xll}
    \hline
    \textbf{Header} & \textbf{Server bawaan} & \textbf{Aplikasi Bab 7} \\
    \hline
    \texttt{Content-Security-Policy} & \dots & \dots \\
    \hline
    \texttt{X-Content-Type-Options} & \dots & \dots \\
    \hline
    \texttt{Referrer-Policy} & \dots & \dots \\
    \hline
    \texttt{Permissions-Policy} & \dots & \dots \\
    \hline
    Header tambahan pilihan sendiri & \dots & \dots \\
    \hline
  \end{tabularx}
\end{table}

\subsection*{Latihan 2: Cookie dan Aturan Lintas Origin}

Latihan ini membuktikan bahwa izin lintas origin ditentukan daftar, bukan
kebetulan, dan bahwa atribut cookie terbaca langsung dari \textit{response}.

\begin{enumerate}
  \item Baca header \texttt{Set-Cookie} dari alamat profil, lalu catat tiap
        atributnya.
  \item Minta alamat yang sama tiga kali dengan header \texttt{Origin} yang
        berbeda, seperti pada Listing \ref{lst:c7-lat2-jalan}, lalu catat
        header izin yang kembali ke Tabel \ref{tab:c7-lat2-hasil}. Tabel
        \ref{tab:c7-lat2-contoh} memperlihatkan contoh yang sudah terisi.
  \item Tambahkan satu origin baru ke daftar di
        \berkas{sources/chapter07/kebijakan.json}, nyalakan ulang
        aplikasinya, lalu ulangi pemeriksaannya.
  \item Jawab tiga pertanyaan berikut. Pertama, siapa yang sebenarnya
        dilindungi aturan CORS? Kedua, mengapa aplikasi ini tidak memakai
        nilai bintang pada header izin? Ketiga, mengapa \texttt{Vary: Origin}
        tetap dikirim walaupun izinnya tidak diberikan?
\end{enumerate}

\begin{lstlisting}[language=bash, caption={Membaca cookie dan izin lintas origin, tersedia di \berkas{sources/chapter07/aplikasi.py}}, label={lst:c7-lat2-jalan}]
curl -sI http://127.0.0.1:8030/api/profil | grep -i "^set-cookie"
curl -sI -H "Origin: https://toko.pradita.ac.id" \
  http://127.0.0.1:8030/api/profil | grep -i "^access-control"
curl -sI -H "Origin: https://situs-lain.example" \
  http://127.0.0.1:8030/api/profil | grep -iE "^(access-control|vary)"

# Keluarannya:
# Set-Cookie: sesi=...; HttpOnly; SameSite=Lax; Path=/; Max-Age=3600
# Access-Control-Allow-Origin: https://toko.pradita.ac.id
# Vary: Origin
\end{lstlisting}

\begin{table}[!htb]
  \centering
  \caption{Contoh hasil Latihan 2 yang sudah terisi}
  \label{tab:c7-lat2-contoh}
  \begin{tabularx}{\textwidth}{lX}
    \hline
    \textbf{Origin yang dikirim} & \textbf{Header izin yang kembali} \\
    \hline
    \texttt{https://toko.pradita.ac.id} & berisi origin itu sendiri \\
    \hline
    \texttt{https://situs-lain.example} & tidak ada, hanya \texttt{Vary: Origin} \\
    \hline
  \end{tabularx}
\end{table}

\begin{table}[!htb]
  \centering
  \caption{Lembar isian hasil pengamatan Latihan 2}
  \label{tab:c7-lat2-hasil}
  \begin{tabularx}{\textwidth}{lX}
    \hline
    \textbf{Yang diperiksa} & \textbf{Hasilnya} \\
    \hline
    Atribut cookie sesi & \dots \\
    \hline
    Origin terdaftar pertama & \dots \\
    \hline
    Origin terdaftar kedua & \dots \\
    \hline
    Origin tidak terdaftar & \dots \\
    \hline
    Origin baru yang ditambahkan sendiri & \dots \\
    \hline
  \end{tabularx}
\end{table}

\subsection*{Latihan 3: Teks Pengguna sebagai Data}

Latihan ini membuktikan bahwa teks bertanda baca tidak membuat basis data
maupun browser berpindah tugas.

\begin{enumerate}
  \item Jalankan pencarian biasa, lalu catat jumlah barisnya.
  \item Jalankan pencarian dengan kata kunci yang di dalamnya ada tanda kutip
        dan tanda titik koma, seperti pada Listing \ref{lst:c7-lat3-jalan},
        lalu catat jumlah baris dan status permintaannya.
  \item Buka halaman sapaan dengan nama yang berisi tanda kurung siku, lalu
        catat apa yang tertulis pada sumber halamannya.
  \item Isi Tabel \ref{tab:c7-lat3-hasil}. Tabel \ref{tab:c7-lat3-contoh}
        memperlihatkan contoh yang sudah terisi.
  \item Jawab dua pertanyaan berikut. Pertama, mengapa pencarian kedua
        menghasilkan nol baris, dan apa yang akan terjadi bila teks itu
        disambung langsung ke SQL? Kedua, apa bedanya menolak teks tertentu
        dengan memperlakukan seluruh teks sebagai data?
\end{enumerate}

\begin{lstlisting}[language=bash, caption={Mengirim teks bertanda baca sebagai kata kunci, tersedia di \berkas{sources/chapter07/aplikasi.py}}, label={lst:c7-lat3-jalan}]
curl -s "http://127.0.0.1:8030/api/cari?judul=Pengantar%20Jilid%207" | head -c 120
curl -s --get --data-urlencode "judul=Pengantar'; SELECT 1 --" \
  http://127.0.0.1:8030/api/cari | head -c 120
curl -s --get --data-urlencode 'nama=<b>Mahasiswa</b>' \
  http://127.0.0.1:8030/sapa | grep -i halo

# Keluarannya:
# {"kata_kunci": "Pengantar Jilid 7", "hasil": [{"id": ...
# {"kata_kunci": "Pengantar'; SELECT 1 --", "hasil": []}
# <body><h1>Halo, &lt;b&gt;Mahasiswa&lt;/b&gt;</h1>
\end{lstlisting}

\begin{table}[!htb]
  \centering
  \caption{Contoh hasil Latihan 3 yang sudah terisi}
  \label{tab:c7-lat3-contoh}
  \begin{tabularx}{\textwidth}{Xrl}
    \hline
    \textbf{Yang dikirim} & \textbf{Baris} & \textbf{Status} \\
    \hline
    \texttt{Pengantar Jilid 7} & 20 & 200 \\
    \hline
    \texttt{Pengantar'; SELECT 1 -{}-} & 0 & 200 \\
    \hline
  \end{tabularx}
\end{table}

\begin{table}[!htb]
  \centering
  \caption{Lembar isian hasil pengamatan Latihan 3}
  \label{tab:c7-lat3-hasil}
  \begin{tabularx}{\textwidth}{Xrl}
    \hline
    \textbf{Yang dikirim} & \textbf{Baris} & \textbf{Status} \\
    \hline
    Kata kunci biasa & \dots & \dots \\
    \hline
    Kata kunci bertanda baca & \dots & \dots \\
    \hline
    Nama berisi tanda kurung siku & \dots & \dots \\
    \hline
  \end{tabularx}
\end{table}

\subsection*{Latihan 4: Menegakkan Gate Keamanan}

Latihan ini memasang seluruh syarat di atas sebagai satu gerbang, lalu
membandingkan vonisnya pada dua sasaran.

\begin{enumerate}
  \item Baca \berkas{sources/chapter07/kebijakan.json}, lalu tulis dugaan
        berapa syarat yang akan terpenuhi aplikasi dan berapa pada server
        pembanding.
  \item Jalankan gerbangnya terhadap keduanya seperti pada Listing
        \ref{lst:c7-lat4-jalan}, lalu catat jumlah syarat, jumlah yang lolos,
        dan status keluarnya ke Tabel \ref{tab:c7-lat4-hasil}. Tabel
        \ref{tab:c7-lat4-contoh} memperlihatkan contoh yang sudah terisi.
  \item Jalankan pemeriksaan dependensi dan analisis statis, lalu catat
        jumlah temuannya.
  \item Hapus sementara satu header dari \berkas{sources/chapter07/kebijakan.json},
        nyalakan ulang aplikasinya, jalankan ulang gerbangnya, lalu
        kembalikan seperti semula.
  \item Jawab tiga pertanyaan berikut. Pertama, mengapa gerbang menilai
        \textit{response} yang dikirim dan bukan membaca kodenya? Kedua,
        mengapa satu syarat justru lolos pada server pembanding? Ketiga,
        mengapa hasil pemeriksaan dependensi punya umur simpan?
\end{enumerate}

\begin{lstlisting}[language=bash, caption={Menjalankan gate pada dua sasaran, tersedia di \berkas{sources/chapter07/periksa-keamanan.py}}, label={lst:c7-lat4-jalan}]
source ../.venv/bin/activate
python periksa-keamanan.py
echo $?
python periksa-keamanan.py kebijakan-pembanding.json
echo $?
./audit-dependensi.sh
ruff check .

# Keluarannya:
# Lolos: seluruh syarat keamanan terpenuhi.
# 0
# Gagal: 8 syarat keamanan belum terpenuhi.
# 1
# No known vulnerabilities found
# All checks passed!
\end{lstlisting}

\begin{table}[!htb]
  \centering
  \caption{Contoh hasil Latihan 4 yang sudah terisi}
  \label{tab:c7-lat4-contoh}
  \begin{tabularx}{\textwidth}{Xrrr}
    \hline
    \textbf{Sasaran} & \textbf{Syarat} & \textbf{Lolos} & \textbf{Status keluar} \\
    \hline
    Aplikasi Bab 7 & 18 & 18 & 0 \\
    \hline
    Server statis bawaan Python & 9 & 1 & 1 \\
    \hline
  \end{tabularx}
\end{table}

\begin{table}[!htb]
  \centering
  \caption{Lembar isian hasil pengamatan Latihan 4}
  \label{tab:c7-lat4-hasil}
  \begin{tabularx}{\textwidth}{Xrrr}
    \hline
    \textbf{Sasaran} & \textbf{Syarat} & \textbf{Lolos} & \textbf{Status keluar} \\
    \hline
    Aplikasi Bab 7 & \dots & \dots & \dots \\
    \hline
    Server statis bawaan Python & \dots & \dots & \dots \\
    \hline
    Aplikasi dengan satu header dihapus & \dots & \dots & \dots \\
    \hline
    Temuan dependensi dan analisis statis & \dots & \dots & \dots \\
    \hline
  \end{tabularx}
\end{table}

Setelah seluruh latihan selesai, container dihentikan dengan
\texttt{docker compose down -v}. Opsi \texttt{-v} ikut menghapus volume
datanya, sehingga penyiapan berikutnya dimulai dari keadaan bersih.
